\documentclass[aspectratio=169,12pt]{beamer}

% ---------------------------------------------------------------
% Slides para projetar as quest\~oes do Bloco 5 (fazer em sala).
% Mesma identidade dos outros documentos de din\^amica:
% Palatino, versalete, uma cor de destaque (bord\^o).
% Os n\'umeros das quest\~oes s\~ao os mesmos da lista impressa.
% ---------------------------------------------------------------
\usepackage[utf8]{inputenc}
\usepackage[T1]{fontenc}
\usepackage[brazil]{babel}
\usepackage[sc]{mathpazo}   % sem 'osf': algarismos alinhados leem melhor projetados
\usefonttheme{serif}
\usepackage{microtype}
\usepackage{amsmath}
\usepackage{amssymb}
\usepackage{xcolor}
\usepackage{tikz}
\usepackage{pgfplots}
\pgfplotsset{compat=1.18}
\usetikzlibrary{arrows.meta, patterns}

\definecolor{vinho}{RGB}{122, 27, 42}
\definecolor{grafite}{RGB}{70, 70, 70}

\setbeamercolor{structure}{fg=vinho}
\setbeamercolor{normal text}{fg=black}
\setbeamercolor{alerted text}{fg=vinho}
\setbeamertemplate{navigation symbols}{}
\setbeamersize{text margin left=0.9cm, text margin right=0.9cm}

% ---------------------------------------------------------------
% T\'itulo de slide: versalete bord\^o \`a esquerda, fonte/tempo \`a
% direita, filete embaixo
% ---------------------------------------------------------------
\setbeamertemplate{frametitle}{%
  \vspace{0.18cm}%
  \begin{minipage}{\textwidth}
    {\large\scshape\color{vinho}\insertframetitle}%
    \hfill{\small\itshape\color{grafite}\insertframesubtitle}\\[-0.45em]
    {\color{vinho}\rule{\textwidth}{1.1pt}}
  \end{minipage}%
  \par\vspace{0.08cm}
}

\setbeamertemplate{footline}{%
  \hfill{\scriptsize\color{black!45}%
    din\^amica \,\textperiodcentered\, leis de newton
    \,\textperiodcentered\, \insertframenumber}%
  \hspace{0.9cm}\vspace{0.25cm}%
}

% ---------------------------------------------------------------
% Alternativas
% ---------------------------------------------------------------
\newcommand{\optitem}[2]{%
  \par\hangindent=1.5em\hangafter=1
  \makebox[1.5em][l]{\textbf{#1)}}#2}

\newcommand{\alternativas}[5]{%
  \par\vspace{0.55em}%
  \begingroup\setlength{\parskip}{0.30em}%
  \optitem{a}{#1}\optitem{b}{#2}\optitem{c}{#3}\optitem{d}{#4}\optitem{e}{#5}%
  \par\endgroup}

% vers\~ao em duas colunas (a, b, c | d, e), para alternativas m\'edias
\newcommand{\alternativascols}[5]{%
  \par\vspace{0.5em}%
  \begingroup\setlength{\parskip}{0.28em}%
  \noindent
  \begin{minipage}[t]{0.48\textwidth}
    \optitem{a}{#1}\optitem{b}{#2}\optitem{c}{#3}
  \end{minipage}\hfill
  \begin{minipage}[t]{0.48\textwidth}
    \optitem{d}{#4}\optitem{e}{#5}
  \end{minipage}%
  \par\endgroup}

% vers\~ao em uma linha s\'o, para alternativas curtas (num\'ericas)
\newcommand{\alternativascurtas}[5]{%
  \par\vspace{0.7em}%
  \noindent\textbf{a)}~#1\hfill\textbf{b)}~#2\hfill\textbf{c)}~#3%
  \hfill\textbf{d)}~#4\hfill\textbf{e)}~#5\par}

% ---------------------------------------------------------------
% Slide de resposta (projetar depois de resolver no quadro)
% ---------------------------------------------------------------
\newcommand{\resposta}[2]{%
  \par\vspace{0.2em}
  {\Large\bfseries\color{vinho}Quest\~ao #1 \,\textemdash\, resposta: #2}
  \par\vspace{0.6em}}

\setlength{\parskip}{0.35em}
\linespread{1.03}

% ---------------------------------------------------------------
\begin{document}
% ---------------------------------------------------------------

% ===================================================================
% CAPA
% ===================================================================
\begin{frame}[plain]
\vspace{0.6cm}
{\scshape\color{vinho}cursinho solid\'ario \,\textperiodcentered\, f\'isica}\\[-0.5em]
{\color{vinho}\rule{\textwidth}{1.6pt}}\\[1.0em]
{\huge\bfseries Quest\~oes em Sala}\\[0.45em]
{\LARGE Din\^amica --- Leis de Newton}\\[1.0em]
{\itshape\color{grafite}Bloco 5 do plano de aula \,\textperiodcentered\,
8 quest\~oes \,\textperiodcentered\, 35 min}
\hfill {\color{grafite}2026}\\[-0.3em]
{\color{black!30}\rule{\textwidth}{0.5pt}}\\[0.8em]
{\small\color{grafite}A numera\c{c}\~ao \'e a mesma da lista impressa
(\texttt{lista\_questoes\_dinamica.pdf}).}
\end{frame}

% ===================================================================
% ROTEIRO
% ===================================================================
\begin{frame}
\frametitle{Roteiro}
\framesubtitle{35 min}

\small
{\renewcommand{\arraystretch}{1.22}
\begin{tabular}{@{}clll@{}}
  \textbf{\#} & \textbf{Fonte} & \textbf{O que cobra} & \textbf{Tempo}\\[0.15em]
  \hline
  \rule{0pt}{1.15em}%
  Q1  & ENEM PPL 2012      & 1\textordfeminine{} Lei (bola no vag\~ao)      & 4 min \\
  Q9  & ENEM PPL 2012      & 3\textordfeminine{} Lei (m\~ae e o m\'ovel)    & 3 min \\
  Q8  & ENEM 2013          & 3\textordfeminine{} Lei + atrito (rampa)       & 4 min \\
  Q4  & UTFPR Ver\~ao 2025 & 2\textordfeminine{} Lei num\'erica             & 6 min \\
  Q5  & ENEM 2017          & 2\textordfeminine{} Lei + gr\'afico (cinto)    & 5 min \\
  Q12 & UEL 1995           & Plano inclinado: $a = g\,\mathrm{sen}\,\theta$ & 4 min \\
  Q13 & UERJ 1998          & Rampa: for\c{c}a m\'inima                      & 4 min \\
  Q14 & UTFPR Inverno 2024 & Blocos em contato                              & 5 min \\
\end{tabular}}

\vspace{0.6em}
{\small\itshape\color{grafite}Para cada uma: ler o enunciado
\,\textperiodcentered\, 2 min tentando \,\textperiodcentered\, votar
nas alternativas \,\textperiodcentered\, resolver no quadro.}
\end{frame}

% ===================================================================
% Q1
% ===================================================================
\begin{frame}
\frametitle{Quest\~ao 1}
\framesubtitle{ENEM PPL 2012 \,\textperiodcentered\, $\sim$4 min}

\footnotesize\linespread{0.97}\selectfont
Em 1543, Nicolau Cop\'ernico publicou um livro revolucion\'ario em que
propunha a Terra girando em torno do seu pr\'oprio eixo e rodando em
torno do Sol. Isso contraria a concep\c{c}\~ao aristot\'elica, que
acredita que a Terra \'e o centro do universo. Para os aristot\'elicos,
se a Terra gira do oeste para o leste, coisas como nuvens e p\'assaros,
que n\~ao est\~ao presas \`a Terra, pareceriam estar sempre se movendo
do leste para o oeste, justamente como o Sol. Mas foi Galileu Galilei
que, em 1632, baseando-se em experi\^encias, rebateu a cr\'itica
aristot\'elica, confirmando assim o sistema de Cop\'ernico. Seu
argumento, adaptado para a nossa \'epoca, \'e: \textbf{se uma pessoa,
dentro de um vag\~ao de trem em repouso, solta uma bola, ela cai junto
a seus p\'es. Mas se o vag\~ao estiver se movendo com velocidade
constante, a bola tamb\'em cai junto a seus p\'es.} Isto porque a bola,
enquanto cai, continua a compartilhar do movimento do vag\~ao.

\vspace{0.2em}
\small
O princ\'ipio f\'isico usado por Galileu para rebater o argumento
aristot\'elico foi

\alternativascols{a lei da in\'ercia.}{a\c{c}\~ao e rea\c{c}\~ao.}%
{a 2\textordfeminine{} lei de Newton.}{a conserva\c{c}\~ao da energia.}%
{o princ\'ipio da equival\^encia.}
\end{frame}

\begin{frame}
\frametitle{Quest\~ao 1}
\framesubtitle{gabarito}
\resposta{1}{A --- lei da in\'ercia}

Enquanto cai, a bola \textbf{mant\'em por in\'ercia} a velocidade
horizontal que tinha junto com o vag\~ao --- a resultante horizontal
nela \'e nula.

\vspace{0.2em}
Por isso ela cai junto aos p\'es. \textbf{Tanto faz} o vag\~ao estar
parado ou em MRU: nos dois casos $\vec{F}_R = \vec{0}$ na horizontal.

\vspace{0.9em}
{\color{grafite}\itshape Pegadinha: quem responde ``2\textordfeminine{}
Lei'' est\'a procurando uma for\c{c}a que n\~ao existe. Sem for\c{c}a
horizontal, o movimento horizontal simplesmente continua.}
\end{frame}

% ===================================================================
% Q9
% ===================================================================
\begin{frame}
\frametitle{Quest\~ao 9}
\framesubtitle{ENEM PPL 2012 \,\textperiodcentered\, $\sim$3 min}

\small
Durante uma faxina, a m\~ae pediu que o filho a ajudasse, deslocando um
m\'ovel para mud\'a-lo de lugar. Para escapar da tarefa, o filho disse
ter aprendido na escola que n\~ao poderia puxar o m\'ovel, pois a
Terceira Lei de Newton define que \textbf{se puxar o m\'ovel, o m\'ovel
o puxar\'a igualmente de volta}, e assim n\~ao conseguir\'a exercer uma
for\c{c}a que possa coloc\'a-lo em movimento.

\vspace{0.3em}
Qual argumento a m\~ae utilizar\'a para apontar o erro de
interpreta\c{c}\~ao do garoto?

\alternativas{A for\c{c}a de a\c{c}\~ao \'e aquela exercida pelo
garoto.}%
{A for\c{c}a resultante sobre o m\'ovel \'e sempre nula.}%
{As for\c{c}as que o ch\~ao exerce sobre o garoto se anulam.}%
{A for\c{c}a de a\c{c}\~ao \'e um pouco maior que a for\c{c}a de
rea\c{c}\~ao.}%
{O par de for\c{c}as de a\c{c}\~ao e rea\c{c}\~ao n\~ao atua em um
mesmo corpo.}
\end{frame}

\begin{frame}
\frametitle{Quest\~ao 9}
\framesubtitle{gabarito}
\resposta{9}{E --- o par n\~ao atua no mesmo corpo}

A\c{c}\~ao e rea\c{c}\~ao t\^em a \textbf{mesma intensidade} e
\textbf{sentidos opostos}, mas atuam em \textbf{corpos diferentes} ---
ent\~ao \textbf{nunca se cancelam}.

\vspace{0.2em}
A resultante \textit{no m\'ovel} pode ser diferente de zero, e ele
acelera.

\vspace{0.9em}
{\color{grafite}\itshape A armadilha cl\'assica do ENEM: ``a\c{c}\~ao e
rea\c{c}\~ao se anulam, logo nada se move''. S\'o se anulam for\c{c}as
que agem \textbf{no mesmo corpo}.}
\end{frame}

% ===================================================================
% Q8
% ===================================================================
\begin{frame}
\frametitle{Quest\~ao 8}
\framesubtitle{ENEM 2013 \,\textperiodcentered\, $\sim$4 min}

Uma pessoa necessita da for\c{c}a de atrito em seus p\'es para se
deslocar sobre uma superf\'icie. Logo, uma pessoa que sobe uma rampa em
linha reta ser\'a auxiliada pela for\c{c}a de atrito exercida pelo
ch\~ao em seus p\'es.

\vspace{0.3em}
Em rela\c{c}\~ao ao movimento dessa pessoa, quais s\~ao a
\textbf{dire\c{c}\~ao} e o \textbf{sentido} da for\c{c}a de atrito
mencionada no texto?

\alternativas{Perpendicular ao plano e no mesmo sentido do movimento.}%
{Paralelo ao plano e no sentido contr\'ario ao movimento.}%
{Paralelo ao plano e no mesmo sentido do movimento.}%
{Horizontal e no mesmo sentido do movimento.}%
{Vertical e sentido para cima.}
\end{frame}

\begin{frame}
\frametitle{Quest\~ao 8}
\framesubtitle{gabarito}
\resposta{8}{C --- paralelo ao plano, no sentido do movimento}

\begin{columns}[T]
\begin{column}{0.56\textwidth}
Ao andar, o p\'e \textbf{empurra o ch\~ao para tr\'as} (paralelo \`a
rampa).

\vspace{0.3em}
Pela 3\textordfeminine{} Lei, o ch\~ao devolve: o \textbf{atrito}
empurra a pessoa \textbf{para a frente}.

\vspace{0.5em}
{\color{grafite}\itshape Quem anda \'e movido pelo ch\~ao. Sem atrito,
a pessoa s\'o patina no lugar.}
\end{column}
\begin{column}{0.42\textwidth}
\centering
\begin{tikzpicture}[scale=0.92]
  % rampa (inclina\c{c}\~ao de 31\textdegree{}: subida 0,6 por 1 de base)
  \draw[thick, fill=black!7] (0,0) -- (4.5,2.7) -- (4.5,0) -- cycle;
  % pessoa, em p\'e sobre a rampa
  \begin{scope}[shift={(1.85,1.30)}, rotate=31]
    \filldraw[black!75] (0,0.78) circle (0.17);
    \draw[very thick, black!75] (0,0.61) -- (0,0.24);
    \draw[very thick, black!75] (0,0.24) -- (-0.30,0.01);
    \draw[very thick, black!75] (0,0.24) -- (0.28,0.01);
  \end{scope}
  % atrito: empurra a pessoa rampa acima
  \draw[-{Stealth[length=3.4mm]}, very thick, vinho]
    (2.05,1.28) -- (3.05,1.88);
  \node[font=\small, vinho, anchor=south west] at (2.95,1.86) {$f_{at}$};
  % rea\c{c}\~ao: o p\'e empurra o ch\~ao rampa abaixo
  \draw[-{Stealth[length=3mm]}, thick, black!55]
    (1.62,1.11) -- (0.72,0.57);
  \node[font=\scriptsize, black!55, align=center, anchor=north]
    at (1.15,-0.12) {o p\'e empurra\\o ch\~ao para tr\'as};
\end{tikzpicture}
\end{column}
\end{columns}
\end{frame}

% ===================================================================
% Q4
% ===================================================================
\begin{frame}
\frametitle{Quest\~ao 4}
\framesubtitle{UTFPR Ver\~ao 2025 \,\textperiodcentered\, $\sim$6 min}

\small
Sobre um piso horizontal, desloca-se um corpo de massa igual a
$4{,}0$ kg, movendo-se \textbf{sem atrito} e com velocidade inicial de
$18{,}0$ km/h, quando \'e submetido a uma for\c{c}a de $16{,}0$ N
durante um intervalo de tempo de $5{,}0$ s.

\vspace{0.4em}
Sobre este movimento, determine respectivamente:

\vspace{0.2em}
\hspace{1em}1) a acelera\c{c}\~ao do corpo;\quad
2) a velocidade ao final dos $5{,}0$ s;\quad
3) o espa\c{c}o percorrido.

\alternativas{4 m/s\textsuperscript{2}, 25 m/s, 75 m.}%
{2 m/s\textsuperscript{2}, 30 m/s, 200 m.}%
{4 m/s\textsuperscript{2}, 15 m/s, 200 m.}%
{3 m/s\textsuperscript{2}, 18 m/s, 100 m.}%
{4 m/s\textsuperscript{2}, 20 m/s, 200 m.}
\end{frame}

\begin{frame}
\frametitle{Quest\~ao 4}
\framesubtitle{gabarito}
\resposta{4}{A}

\begin{columns}[T]
\begin{column}{0.46\textwidth}
\textbf{1) 2\textordfeminine{} Lei}
\[ a = \frac{F}{m} = \frac{16}{4} = 4\ \text{m/s}^2 \]

\textbf{2) Convers\~ao primeiro!}
\[ v_0 = \frac{18}{3{,}6} = 5\ \text{m/s} \]
\[ v = v_0 + at = 5 + 4\cdot 5 = 25\ \text{m/s} \]
\end{column}
\begin{column}{0.52\textwidth}
\textbf{3) Cinem\'atica}
\[ \Delta s = v_0 t + \tfrac{1}{2}a t^2 \]
\[ \Delta s = 5\cdot 5 + \tfrac{1}{2}\cdot 4\cdot 25 = 75\ \text{m} \]

\vspace{0.6em}
{\color{grafite}\itshape Estilo UTFPR: a f\'isica \'e s\'o o
$a = F/m$; o resto \'e a cinem\'atica da aula passada. Quem esquece de
passar km/h para m/s cai na alternativa errada.}
\end{column}
\end{columns}
\end{frame}

% ===================================================================
% Q5
% ===================================================================
\begin{frame}
\frametitle{Quest\~ao 5}
\framesubtitle{ENEM 2017 \,\textperiodcentered\, $\sim$5 min}

\footnotesize
Em uma colis\~ao frontal entre dois autom\'oveis, a for\c{c}a que o cinto de
seguran\c{c}a exerce sobre o t\'orax e abd\^omen do motorista pode causar
les\~oes graves nos \'org\~aos internos. Um fabricante realizou testes em cinco
modelos diferentes de cinto, simulando uma colis\~ao de $0{,}30$ s de
dura\c{c}\~ao, com bonecos equipados com aceler\^ometros.
\textbf{A massa dos bonecos, as dimens\~oes dos cintos e a velocidade antes e
depois do impacto foram as mesmas em todos os testes.}

\vspace{-0.1em}
\begin{columns}[T]
\begin{column}{0.42\textwidth}
\vspace{0.9em}
\normalsize
Qual modelo de cinto oferece \textbf{menor risco de les\~ao interna} ao
motorista?

\alternativascurtas{1}{2}{3}{4}{5}
\end{column}
\begin{column}{0.56\textwidth}
\centering
\begin{tikzpicture}
\begin{axis}[
  width=7.9cm, height=3.95cm,
  xlabel={tempo (s)}, ylabel={desacelera\c{c}\~ao},
  label style={font=\small},
  tick label style={font=\small},
  xmin=0, xmax=0.32, ymin=0, ymax=110,
  ytick=\empty,
  xtick={0,0.10,0.20,0.30},
  legend pos=north east,
  legend style={font=\scriptsize, draw=black!30},
  axis lines=left,
]
\addplot[thick, smooth, black] coordinates
  {(0,0) (0.03,60) (0.06,100) (0.09,55) (0.13,15) (0.20,5) (0.30,0)};
\addlegendentry{1}
\addplot[very thick, smooth, vinho, dashed] coordinates
  {(0,0) (0.05,30) (0.10,38) (0.15,40) (0.20,38) (0.25,30) (0.30,0)};
\addlegendentry{2}
\addplot[thick, smooth, black!55] coordinates
  {(0,0) (0.07,25) (0.12,80) (0.16,50) (0.22,20) (0.30,0)};
\addlegendentry{3}
\addplot[thick, smooth, black, dotted] coordinates
  {(0,0) (0.10,15) (0.17,30) (0.22,90) (0.26,40) (0.30,0)};
\addlegendentry{4}
\addplot[thick, smooth, black!55, dashdotted] coordinates
  {(0,0) (0.04,45) (0.09,70) (0.15,65) (0.21,25) (0.30,0)};
\addlegendentry{5}
\end{axis}
\end{tikzpicture}\\[-0.2em]
{\scriptsize\itshape (figura redesenhada)}
\end{column}
\end{columns}
\end{frame}

\begin{frame}
\frametitle{Quest\~ao 5}
\framesubtitle{gabarito}
\resposta{5}{B --- o cinto 2}

Mesma massa e mesma varia\c{c}\~ao de velocidade em todos os testes
$\Rightarrow$ a \textbf{\'area embaixo de cada curva \'e igual}.

\vspace{0.3em}
Por $F = ma$, o cinto mais seguro \'e o que tem o \textbf{menor pico de
desacelera\c{c}\~ao} --- ou seja, o que \textbf{espalha a frenagem ao
longo do tempo}. \'E a curva 2, baixa e larga.

\vspace{0.8em}
{\color{grafite}\itshape A mesma ideia do airbag, do capacete e da
zona de deforma\c{c}\~ao do carro: n\~ao d\'a para mudar o
$\Delta v$, ent\~ao aumenta-se o $\Delta t$ para derrubar a for\c{c}a.}
\end{frame}

% ===================================================================
% Q12
% ===================================================================
\begin{frame}
\frametitle{Quest\~ao 12}
\framesubtitle{UEL 1995 \,\textperiodcentered\, $\sim$4 min}

\begin{columns}[T]
\begin{column}{0.56\textwidth}
Um corpo de massa $2{,}0$ kg \'e abandonado sobre um plano
\textbf{perfeitamente liso} e inclinado de $37^{\circ}$ com a
horizontal.

\vspace{0.4em}
Adotando $g = 10$ m/s\textsuperscript{2},
$\mathrm{sen}\,37^{\circ} = 0{,}60$ e $\cos 37^{\circ} = 0{,}80$,
conclui-se que a acelera\c{c}\~ao com que o corpo desce o plano tem
m\'odulo, em m/s\textsuperscript{2}:

\alternativascurtas{4,0}{5,0}{6,0}{8,0}{10}
\end{column}
\begin{column}{0.42\textwidth}
\centering
\begin{tikzpicture}[scale=1.0]
  \draw[thick] (0,0) -- (4.3,0);
  \draw[thick] (0,0) -- (4.3,3.24) -- (4.3,0);
  \draw (0.9,0) arc (0:37:0.9);
  \node at (1.18,0.26) {\small $37^{\circ}$};
  \begin{scope}[shift={(2.1,1.58)}, rotate=37]
    \draw[thick, fill=black!8] (0,0) rectangle (0.95,0.6);
    \node[font=\small] at (0.48,0.3) {2 kg};
  \end{scope}
\end{tikzpicture}
\end{column}
\end{columns}
\end{frame}

\begin{frame}
\frametitle{Quest\~ao 12}
\framesubtitle{gabarito}
\resposta{12}{C --- $6{,}0$ m/s\textsuperscript{2}}

\begin{columns}[T]
\begin{column}{0.54\textwidth}
Na rampa, decompor o \textbf{peso}:
\[ P_x = mg\,\mathrm{sen}\,\theta \qquad
   P_y = mg\cos\theta = N \]

Sem atrito, a resultante \'e s\'o $P_x$:
\[ ma = mg\,\mathrm{sen}\,\theta \;\Rightarrow\;
   a = g\,\mathrm{sen}\,\theta \]
\[ a = 10 \cdot 0{,}60 = 6{,}0\ \text{m/s}^2 \]
\end{column}
\begin{column}{0.44\textwidth}
\centering
\begin{tikzpicture}[scale=0.95]
  \draw[thick] (0,0) -- (4.0,0);
  \draw[thick] (0,0) -- (4.0,3.0) -- (4.0,0);
  \begin{scope}[shift={(1.9,1.42)}, rotate=37]
    \draw[thick, fill=black!8] (0,0) rectangle (0.9,0.55);
    \draw[-{Stealth[length=2.6mm]}, thick, vinho] (0.45,0.28) -- (-0.55,0.28)
      node[left, font=\small] {$P_x$};
    \draw[-{Stealth[length=2.6mm]}, thick, black!60] (0.45,0.28) -- (0.45,1.15)
      node[above, font=\small] {$N$};
  \end{scope}
  \draw[-{Stealth[length=2.6mm]}, thick, black!60] (2.35,1.70) -- (2.35,0.55)
    node[midway, right, font=\small] {$P$};
\end{tikzpicture}
\end{column}
\end{columns}

\vspace{0.3em}
{\color{grafite}\itshape A massa cancela: qualquer corpo desceria essa
rampa com a mesma acelera\c{c}\~ao.}
\end{frame}

% ===================================================================
% Q13
% ===================================================================
\begin{frame}
\frametitle{Quest\~ao 13}
\framesubtitle{UERJ 1998 \,\textperiodcentered\, $\sim$4 min}

\begin{columns}[T]
\begin{column}{0.50\textwidth}
\small
O carregador deseja levar um bloco de \textbf{400 N} de peso at\'e a
carroceria do caminh\~ao, a uma altura de $1{,}5$ m, utilizando-se de
um plano inclinado de $3{,}0$ m de comprimento, conforme a figura.

\vspace{0.4em}
\textbf{Desprezando o atrito}, a for\c{c}a m\'inima com que o
carregador deve puxar o bloco, enquanto este sobe a rampa, ser\'a, em
N, de:

\alternativascurtas{100}{150}{200}{400}{500}
\end{column}
\begin{column}{0.48\textwidth}
\centering
\begin{tikzpicture}[scale=1.0]
  \draw[thick] (-0.8,0) -- (5.6,0);
  \draw[thick] (0,0) -- (2.6,1.5) -- (2.6,0);
  \draw[thick] (2.6,1.5) -- (4.9,1.5);
  \draw[thick] (4.9,1.5) -- (4.9,0);
  \node[anchor=south, font=\small] at (3.75,1.55) {carroceria};
  \begin{scope}[shift={(0.95,0.55)}, rotate=30]
    \draw[thick, fill=black!8] (0,0) rectangle (0.9,0.55);
  \end{scope}
  \draw[<->] (-0.18,0.30) -- (2.42,1.80)
    node[midway, above, sloped, font=\small] {$3{,}0$ m};
  \draw[<->] (3.05,0) -- (3.05,1.5)
    node[midway, right, font=\small] {$1{,}5$ m};
\end{tikzpicture}\\[0.3em]
{\scriptsize\itshape (figura redesenhada)}
\end{column}
\end{columns}
\end{frame}

\begin{frame}
\frametitle{Quest\~ao 13}
\framesubtitle{gabarito}
\resposta{13}{C --- 200 N}

Para subir com velocidade constante (for\c{c}a \textbf{m\'inima}),
basta equilibrar a componente do peso ao longo da rampa:
\[ F = P\,\mathrm{sen}\,\theta = P\cdot\frac{h}{L}
     = 400\cdot\frac{1{,}5}{3{,}0} = 200\ \text{N} \]

\vspace{0.4em}
{\color{grafite}\itshape N\~ao precisa nem do \^angulo: $\mathrm{sen}\,
\theta$ \'e o cateto oposto sobre a hipotenusa, que \'e $h/L$.}

\vspace{0.5em}
A rampa como \textbf{m\'aquina simples}: metade da for\c{c}a, mas o
dobro do caminho ($1{,}5$ m de altura contra $3{,}0$ m de rampa).
\end{frame}

% ===================================================================
% Q14
% ===================================================================
\begin{frame}
\frametitle{Quest\~ao 14}
\framesubtitle{UTFPR Inverno 2024 \,\textperiodcentered\, $\sim$5 min}

\begin{columns}[T]
\begin{column}{0.50\textwidth}
Dois blocos encostados um ao outro, em uma superf\'icie plana e
\textbf{sem atrito}, possuem massas de, respectivamente, 5 kg e 2 kg.
Uma for\c{c}a $F = 140$ N empurra os blocos da esquerda para a direita,
como representado na figura.

\vspace{0.4em}
Qual \'e a \textbf{acelera\c{c}\~ao do sistema} e a \textbf{for\c{c}a
normal entre os dois blocos}, respectivamente?
($g = 10$ m/s\textsuperscript{2}.)
\end{column}
\begin{column}{0.48\textwidth}
\centering
\begin{tikzpicture}[scale=1.05]
  \draw[thick] (-0.2,0) -- (5.2,0);
  \draw[thick, fill=black!8] (1.4,0) rectangle (2.9,1.2);
  \node at (2.15,0.6) {5 kg};
  \draw[thick, fill=black!8] (2.9,0) rectangle (4.0,0.85);
  \node at (3.45,0.42) {2 kg};
  \draw[-{Stealth[length=3mm]}, very thick] (0.15,0.6) -- (1.3,0.6)
    node[midway, above] {$F$};
\end{tikzpicture}\\[0.3em]
{\scriptsize\itshape (figura redesenhada)}
\end{column}
\end{columns}

\vspace{-0.2em}
\alternativascols{10 m/s\textsuperscript{2} e 40 N}%
{10 m/s\textsuperscript{2} e 20 N}%
{10 m/s\textsuperscript{2} e 80 N}%
{20 m/s\textsuperscript{2} e 20 N}%
{20 m/s\textsuperscript{2} e 40 N}
\end{frame}

\begin{frame}
\frametitle{Quest\~ao 14}
\framesubtitle{gabarito}
\resposta{14}{E --- 20 m/s\textsuperscript{2} e 40 N}

\begin{columns}[T]
\begin{column}{0.48\textwidth}
\textbf{Passo 1 --- sistema inteiro}
\[ a = \frac{F}{m_1+m_2} = \frac{140}{7} = 20\ \text{m/s}^2 \]

\textbf{Passo 2 --- s\'o o bloco de 2 kg}\\
A \'unica for\c{c}a horizontal nele \'e a de contato:
\[ N = m_2\,a = 2\cdot 20 = 40\ \text{N} \]
\end{column}
\begin{column}{0.48\textwidth}
\centering
\begin{tikzpicture}[scale=1.0]
  \draw[thick] (0,0) -- (2.2,0);
  \draw[thick, fill=black!8] (0.6,0) rectangle (1.7,0.85);
  \node[font=\small] at (1.15,0.42) {2 kg};
  \draw[-{Stealth[length=2.6mm]}, very thick, vinho] (-0.1,0.42) -- (0.55,0.42)
    node[midway, above, font=\small] {$N$};
  \node[font=\small, anchor=north] at (1.15,-0.15) {diagrama isolado};
\end{tikzpicture}

\vspace{0.5em}
{\color{grafite}\itshape\small O m\'etodo geral de blocos: sistema
inteiro primeiro, depois o diagrama de \textbf{um corpo s\'o}.}
\end{column}
\end{columns}
\end{frame}

% ===================================================================
% FECHAMENTO
% ===================================================================
\begin{frame}
\frametitle{As tr\^es leis, em uma tela}
\framesubtitle{fechamento}

\begin{columns}[T]
\begin{column}{0.5\textwidth}
\textbf{1\textordfeminine{} Lei} --- se $\vec{F}_R = \vec{0}$, o estado
de movimento n\~ao muda.
\[ \vec{F}_R = \vec{0} \;\Rightarrow\; \vec{a} = \vec{0} \]

\textbf{2\textordfeminine{} Lei}
\[ \boxed{\vec{F}_R = m\,\vec{a}} \]

\textbf{3\textordfeminine{} Lei} --- A empurra B com $\vec{F}$, B
empurra A com $-\vec{F}$, \textbf{em corpos diferentes}.
\end{column}
\begin{column}{0.46\textwidth}
\textbf{As for\c{c}as da aula}
\[ P = mg \qquad f_{at} = \mu N \]
\[ N \perp \text{superf\'icie} \qquad T \parallel \text{fio} \]

\textbf{Na rampa}
\[ P_x = mg\,\mathrm{sen}\,\theta \qquad
   N = mg\cos\theta \]

\vspace{0.5em}
{\color{grafite}\itshape Toda quest\~ao de hoje saiu daqui. A
1\textordfeminine{} Lei \'e s\'o o caso $a = 0$ da 2\textordfeminine{}.}
\end{column}
\end{columns}
\end{frame}

\end{document}
